\section*{Appendix}

\textbf{Note:} LLMs like DeepSeek and ChatGPT were used for spell check and debugging of the code. 
They also assisted with creating docstrings. 

\section{Code Listings}
\subsection{Negative loss likelihood} \label{appendix:nll}

\begin{lstlisting}[language=Python]
    def loss_nll(y_true: tf.Tensor, y_pred: tf.Tensor) -> tf.Tensor:
        """Computes the negative log-likelihood (NLL)
        
        Parameters
        ----------
        y_true : tf.Tensor
            True target values. Shape: (batch_size,) or (batch_size, 1).
        y_pred : tf.Tensor
            Predicted mean and variance for each input, as output by the network.
            Shape: (batch_size, 2), where y_pred[:, 0] = mean, y_pred[:, 1] = variance.
        
        Returns
        -------
        tf.Tensor
            Scalar tensor representing the mean NLL loss over the batch.
        """
        y_true = tf.squeeze(y_true)  # shape (batch_size,)
        mu = y_pred[:, 0]            # predicted means
        var = y_pred[:, 1]           # predicted variances (must be positive)
        jitter = 1e-8                # numerical stability
        nll = 0.5 * tf.math.log(var + jitter) + 0.5 * tf.square(y_true - mu) / (var + jitter)
        return tf.reduce_mean(nll)
\end{lstlisting}